%% CNN_CPC manuscript draft for Computer Physics Communications.
%% Updated 2026-08-17: full B0--B33 development history and completed PV1 result.
\documentclass[preprint,12pt]{elsarticle}

\usepackage{amsmath,amssymb}
\usepackage{booktabs}
\usepackage{array}
\usepackage{multirow}
\usepackage{enumitem}
\usepackage{url}
\usepackage{hyperref}
\usepackage{xcolor}
\usepackage{longtable}
\usepackage{tabularx}
\usepackage{graphicx}

\newcommand{\code}[1]{\texttt{#1}}
\newcommand{\R}{\mathbb{R}}
\newcommand{\LNzero}{\mathrm{LN}_0}

\journal{Computer Physics Communications}

\begin{document}
\begin{frontmatter}

\title{CNN\_CPC: Leakage-aware weak supervision and hierarchical multi-series knee MRI classification}

\author[a]{Mohammed H. Talafha\corref{cor1}}
\cortext[cor1]{Corresponding author.}
\address[a]{Research Institute of Sciences and Engineering, University of Sharjah, Sharjah, United Arab Emirates}
% TODO before submission: add corresponding-author e-mail.

\begin{abstract}
Knee magnetic-resonance imaging (MRI) classification is challenging when trusted structured labels are scarce, radiology reports are abundant but incomplete, and each examination contains a heterogeneous collection of pulse sequences. We present \code{CNN\_CPC}, a PyTorch framework for leakage-aware, weakly supervised, study-level classification of twelve knee abnormalities. The released development surface contains 4,407 studies and 24,371 MRI series. Fifty-eight studies carry complete expert labels, whereas the remaining 4,349 are primarily accompanied by reports. The frozen B6 report parser maps each report/target cell to positive, negated, uncertain, or unmentioned; report silence is never converted automatically to a negative label. The active B20 model is trained on 3,120 report-supervised studies containing 17,475 eligible MRI series and 14,123 usable target cells. Each real MRI series is represented by sixteen distributed 2.5D slice triplets, a frozen ConvNeXt-Tiny encoder, learned attention pooling to one series token, a study-level Transformer, and twelve learned pathology queries. B20 additionally uses a deterministic post-resize centered 90\% crop.

The software was developed through a controlled B0--B33 experiment ladder. B10--B17 progressively replaced fixed-stream processing with all-series hierarchical aggregation and improved the visual representation, culminating in the frozen B16 report-aligned encoder and the B17 frozen-encoder downstream regime. B18--B22 studied checkpoint selection and spatial focusing; B20 became the active reference after removing the synthetic vignette artifact introduced by B19. B23--B26 investigated supervision density and class coverage but did not justify replacement of the frozen B6 training contract. B27--B33 then tested increasingly constrained complementary series summaries. On the repeatedly reused 58-study expert development surface, B20 replay reached macro ROC AUC 0.66741, B29 reached 0.67689, B31 reached 0.68228, B32 reached 0.66870, and the simplified B33 reached 0.67645. These values are development diagnostics, not independent validation.

To stop further optimization against the reused expert set, prospective weak-validation v1 (PV1) was frozen before B34. A deterministic StudyInstanceUID-only hash split assigns 2,496 B6-active studies to matched training and 624 studies, comprising 3,544 MRI series, to downstream validation. B20, B31, and B33 were retrained under identical fixed-E2 conditions with the same frozen B16 encoder. The predeclared primary metric is the macro average of per-target B6-weighted soft-label binary cross-entropy, where lower is better. PV1 yielded 0.61558 for B20, 0.57431 for B31, and 0.58497 for B33. The paired bootstrap median B31--B20 difference was $-0.04114$ with 95\% interval $[-0.05182,-0.02959]$; B33--B20 was $-0.03069$ with interval $[-0.04735,-0.01303]$; and B33--B31 was $+0.01086$ with interval $[+0.00225,+0.01951]$. The secondary macro AUC was 0.57276, 0.75673, and 0.75652 for B20, B31, and B33, respectively. Thus B31 is the prospectively selected downstream architecture under the frozen PV1 criterion, while B20 remains the active historical model because PV1 uses weak labels and is not independent expert validation. The framework explicitly separates model development, weak-label architecture selection, reused expert diagnostics, and future independent hidden or external evaluation.

\vspace{1em}
\noindent\textbf{PROGRAM SUMMARY}

\begin{small}
\noindent
{\em Program Title:} CNN\_CPC \\
{\em Repository:} \url{https://github.com/mtalafha90/CNN_CPC} \\
{\em Programming language:} Python $\geq 3.10$ with PyTorch. \\
{\em Repository snapshot described here:} 17 August 2026. \\
{\em Active historical model:} B20 (\code{B20\_crop\_only\_joint\_focus}), canonical epoch 2. \\
{\em Prospectively selected downstream development architecture:} B31 under PV1; not promoted to independent-test status. \\

{\em Nature of problem:} Study-level multi-label classification of twelve knee MRI abnormalities from heterogeneous DICOM examinations with sparse trusted labels, incomplete report-derived supervision, variable numbers of acquisitions, protocol heterogeneity, and occasional incomplete DICOM metadata. \\

{\em Solution method:} CNN\_CPC validates metadata, repairs missing series descriptors from DICOM headers when needed, orders slices spatially, constructs distributed 2.5D triplets, encodes them with ConvNeXt-Tiny, pools sixteen slice tokens to one series token, integrates all eligible real series with a study-level Transformer, and predicts twelve pathologies with learned queries. Reports are used only to create training supervision or representation-alignment targets; final prediction is MRI-only. \\

{\em Restrictions and unusual features:} Expert-labelled studies are excluded from weak-training gradients. The 58-study expert surface is explicitly treated as repeatedly reused development evidence. PV1 is a prospective weak-label architecture-selection surface and is valid only for downstream comparisons that share the exact frozen B16 encoder. CNN\_CPC is research software and is not a clinical diagnostic device.
\end{small}
\end{abstract}

\begin{keyword}
knee MRI \sep weak supervision \sep DICOM \sep ConvNeXt \sep Transformer \sep multi-series MRI \sep hierarchical aggregation \sep report supervision \sep prospective validation \sep reproducibility
\end{keyword}

\end{frontmatter}

\section{Introduction}
\label{sec:introduction}

MRI provides complementary information about knee ligaments, menisci, cartilage, bone, synovium, and joint fluid through multiple anatomical planes and pulse sequences. Examination-level knee MRI classification has been demonstrated by systems such as MRNet \cite{bien2018mrnet}, while modern convolutional encoders such as ConvNeXt \cite{liu2022convnext} and attention-based sequence models \cite{vaswani2017attention} provide flexible building blocks for multi-series learning. In the present setting, however, the central difficulty is not simply network capacity. It is how to learn twelve study-level targets when structured labels are scarce, free-text reports are incomplete, acquisition protocols vary, and a very small expert-labelled surface must be protected from both gradient leakage and excessive model-selection reuse.

CNN\_CPC was developed around six rules. Reports are training supervision rather than inference inputs. Report silence is not a negative label. Every eligible real MRI series may contribute evidence. Expert-labelled studies do not enter weak-training gradients. Each model intervention is versioned and frozen before its designated evaluation. Finally, every evaluation surface is assigned an explicit role: the historical 58-study expert set is reused development evidence, PV1 is prospective weak-label architecture-selection evidence, and hidden competition or external expert data are required for an independent performance claim.

The repository contains the complete B0--B33 history rather than only the early B0--B9 experiments. The present manuscript therefore documents the full transition from early fixed-stream baselines, through variable-series hierarchical aggregation and report-aligned representation learning, to the active B20 model, the B23--B26 supervision studies, the B27--B33 complementary-series family, and the completed prospective PV1 comparison.

\section{Problem definition and data contract}
\label{sec:problem}

\subsection{Targets}
For study $i$, the model predicts $\hat{\mathbf y}_i\in[0,1]^{12}$ in the fixed order ACL, MCL, medial meniscus, lateral meniscus, medial osteoarthritis, lateral osteoarthritis, patellofemoral osteoarthritis, effusion, synovitis, Baker's cyst, contusion, and fracture.

\subsection{Released data surface}
The verified release contains 4,407 training studies, 58 fully labelled expert studies, 4,349 non-gold report studies, and 24,371 series rows. The active B20-family weak-supervision surface contains 3,120 studies with at least one usable B6 cell. Across those studies, 17,475 MRI series satisfy the frozen variable-series policy and 14,123 cells are supervised directly: 6,871 positive and 7,252 negated.

\subsection{Evaluation roles}
The 58 expert studies were used repeatedly in historical model selection and checkpoint diagnostics. They are therefore called the \emph{reused expert development surface}, not an independent validation cohort. The historical 623-study weak-v2 subset is also not a valid holdout for full-surface B20-family experiments because those experiments use all 3,120 B6-active studies.

PV1 was frozen before any B34 experiment. It uses only StudyInstanceUID to define the split. PV1 is fresh with respect to supervised downstream B20/B31/B33 architecture selection, but it is not independent clinical validation because labels are derived from B6 and the historical B16 representation stage had already seen the non-gold reports, including reports corresponding to PV1 validation studies.

\section{DICOM ingestion and 2.5D preprocessing}
\label{sec:dicom}

DICOM objects are decoded with pydicom. When \code{ImageOrientationPatient} and \code{ImagePositionPatient} are present, ordering is determined from the image-plane normal
\begin{equation}
\mathbf n=\mathbf r\times\mathbf c,
\end{equation}
and scalar coordinate $z=\mathbf p\cdot\mathbf n$. \code{InstanceNumber} is the fallback. Multi-frame objects are expanded deterministically. Pixel values are rescaled with DICOM slope/intercept, MONOCHROME1 is inverted, robust finite percentiles are used for clipping, and intensities are mapped to the neural preprocessing range.

For each eligible series, sixteen positions are distributed through the volume. With center $c_s$ and triplet gap $g$, the 2.5D input is
\begin{equation}
X_s=[I_{c_s-g},I_{c_s},I_{c_s+g}].
\end{equation}
Each triplet is represented as three channels at $224\times224$. Evaluation uses the frozen center-offset TTA $[-1,0,1]$.

B20 applies its joint-focus operation \emph{after} the historical resize: a deterministic centered crop retaining 90\% of the resized field is taken and resized back to $224\times224$. No vignette, synthetic border, or crop jitter is used.

\section{Weak supervision from radiology reports}
\label{sec:reports}

The frozen B6 v1.2.1 parser assigns each report/target cell
\begin{equation}
r_{ic}\in\{\mathrm{positive},\mathrm{negated},\mathrm{uncertain},\mathrm{unmentioned}\}.
\end{equation}
The downstream mapping is
\begin{center}
\begin{tabular}{lcc}
\toprule
B6 state & Soft target & Base weight \\
\midrule
positive & 0.85 & 0.50 \\
negated & 0.05 & 1.00 \\
uncertain & ignored & 0.00 \\
unmentioned & ignored & 0.00 \\
\bottomrule
\end{tabular}
\end{center}
Target-wise multipliers equalize total supervision mass. The asymmetry reflects the higher reliability of explicit negation relative to positive report mentions in this frozen rule set. Uncertain and unmentioned cells never become implicit negatives.

\section{Hierarchical all-series architecture}
\label{sec:architecture}

For each real MRI series, the visual encoder maps sixteen 2.5D triplets to $X\in\R^{16\times768}$. Historical multi-head attention pooling maps these slice tokens to one series representation $A\in\R^{768}$. The variable set of series representations is passed to the study-level Transformer. Twelve learned pathology queries then produce twelve logits.

The B20 computation is summarized as
\begin{equation}
\text{study}\rightarrow\{\text{eligible series}\}
\rightarrow\{16\text{ triplets/series}\}
\rightarrow\text{ConvNeXt-Tiny}
\rightarrow A_{\mathrm{series}}
\rightarrow\text{study Transformer}
\rightarrow\text{12 pathology queries}.
\end{equation}

The visual encoder is the historical B16 full-report-aligned ConvNeXt-Tiny encoder. It remains frozen in the B17--B33 downstream family. Its canonical SHA-256 fingerprint is
\begin{verbatim}
b328667cf9dfa9b909ef181c1bcc8975ec42bcd8b9eddad08f908875b73fae96
\end{verbatim}
The canonical active B20 checkpoint is \code{runs/b20\_crop\_focus/b20\_model.pt}, epoch 2.

\section{Development governance}
\label{sec:governance}

The development policy separates optimization evidence from independent evidence. The reused 58-study macro AUC may diagnose a frozen experiment but may not justify post-hoc target switching, per-target model mixtures, candidate-specific blend weights, or a ``.1'' model defined after observing target-wise outcomes. Failed formulations are closed rather than locally tuned against that surface.

PV1 may guide future downstream architecture selection only when the exact frozen B16 encoder and frozen split are retained. It does not authorize encoder selection, report-parser tuning, or claims of expert-label generalization. B20 therefore remains the active historical checkpoint even when a later architecture wins PV1.

\section{Development history: B0--B9 foundations}
\label{sec:b0b9}

The first phase established the representation and supervision foundations. Table~\ref{tab:b0b9} records the historical global point estimates on the reused expert surface.

\begin{table}[ht]
\centering
\caption{B0--B9 development ladder. Values are reused-development macro ROC AUCs, not independent validation.}
\label{tab:b0b9}
\begin{tabular}{llc}
\toprule
ID & Main intervention & Macro AUC \\
\midrule
B0 & Random-init Stage-1 & 0.476254 \\
B1 & Competition-only MRI SSL & 0.503028 \\
B2 & Lower encoder learning rate & 0.499324 \\
B3 & Pathology-aware low-capacity MIL & 0.494465 \\
B4 & Frozen SSL features + PCA/LR & 0.513757 \\
B5 & Image--report SSL + frozen probe & 0.524365 \\
B7-v1 & Direct B6 weak supervision & 0.539772 \\
B7.1 & Complete epoch coverage & 0.564480 \\
B8 & $2\times2$ spatial anatomy tokens & 0.530096 \\
B9 & Strict semantic routing & 0.533496 \\
\bottomrule
\end{tabular}
\end{table}

B5 established that image--report representation alignment was useful. B6 then froze pathology-specific report states, and B7/B7.1 showed that direct report-derived pathology supervision improved the MRI learner. B8 and B9 were negative experiments: greater spatial or semantic structure did not improve the global result. These failures motivated a shift away from increasingly rigid fixed-stream routing.

\section{Development history: B10--B17 representation and aggregation}
\label{sec:b10b17}

\subsection{B10 physical-scale normalization}
B10 tested a plane-specific physical-field normalization using DICOM pixel spacing before resizing. It preserved the B7.1 model and B6 supervision. B10 reached macro AUC 0.552398, below B7.1 at 0.564480, so the global physical-normalization replacement was rejected. This did not rule out physical-scale information in general; it closed only this fixed $224\times224$ normalization protocol.

\subsection{B11/B11.1 pseudo-label completion}
B11 attempted to fill B6-silent cells with high-confidence teacher predictions. A fixed absolute confidence rule failed its viability gate. B11.1 replaced the absolute threshold with target-wise 5th/95th percentile teacher tails and added 3,656 pseudo cells while preserving all 14,123 B6 cells. Despite successful training, B11.1 reached 0.550690 on the reused expert surface, below B7.1. The teacher-pseudo-label completion family was therefore closed.

\subsection{B12 variable-number-of-series model}
B12 replaced the six fixed semantic slots with all eligible repaired sagittal, coronal, and axial MRI series. The frozen audit retained 17,475 series across the 3,120 B6-active studies, compared with 15,468 unique series in the historical dual-stream mapping. B12 reached 0.566092, demonstrating that all-series processing was viable and establishing the acquisition representation inherited by later models.

\subsection{B13 hierarchical series tokens and ImageNet ConvNeXt}
B13 paired the all-series hierarchy with ImageNet-pretrained ConvNeXt-Tiny. Sixteen 2.5D positions were encoded per series; an eight-head attention pool compressed each series to one token; a two-layer study Transformer and pathology queries produced the outputs. B13 reached 0.629357. Its paired interval versus B12 was entirely above zero in the historical analysis, making B13 the first large step in the later architecture lineage.

\subsection{B14 full slice-token memory}
B14 tested whether retaining all $K\times16$ slice tokens at the study level would outperform one-token-per-series compression. It reached 0.619791, below B13, despite fitting B6 training loss more strongly. This supplied early evidence that additional downstream capacity and tighter weak-label fitting were not sufficient for improved expert ranking.

\subsection{B15 knee-MRI domain adaptation}
B15 inserted same-study knee-MRI contrastive adaptation between ImageNet initialization and the B13 hierarchy. On the historical weak-v2 teacher-agreement surface it improved macro AUC from 0.565250 to 0.731906, but on the reused expert surface it reached only 0.620900. The large teacher-agreement gain therefore did not transfer to a global expert-set improvement.

\subsection{B16 full-report semantic alignment}
B16 used the completed B15 encoder and aligned MRI representations with full normalized reports represented by TF--IDF, truncated SVD, and L2 normalization. The representation stage used all 4,349 non-gold report studies and excluded the 58 expert studies. Returning to the full B13/B6 downstream surface yielded 0.634977. The gain over B13 was small and statistically unresolved, but the predeclared point-estimate rule retained B16 as the representation source. The report head is training-only; inference remains MRI-only.

\subsection{B17 frozen B16 encoder}
B17 froze the B16 encoder completely and trained only the hierarchy/head for five fixed epochs. The encoder hash remained unchanged. B17 reached 0.642589 at epoch 5, modestly above B16. This supported the working hypothesis that preserving report-aligned visual features may be preferable to continued direct adaptation of the encoder to sparse/noisy B6 supervision.

\begin{table}[ht]
\centering
\caption{B10--B17 historical development summary.}
\label{tab:b10b17}
\begin{tabular}{llc}
\toprule
Model & Main intervention & Reused macro AUC \\
\midrule
B10 & Physical-scale normalization & 0.552398 \\
B11.1 & Quantile teacher tails & 0.550690 \\
B12 & All real MRI series & 0.566092 \\
B13 & Hierarchical series tokens + ImageNet & 0.629357 \\
B14 & Full slice-token memory & 0.619791 \\
B15 & Knee-MRI same-study SSL & 0.620900 \\
B16 & Full-report semantic alignment & 0.634977 \\
B17 & Frozen B16 encoder & 0.642589 \\
\bottomrule
\end{tabular}
\end{table}

\section{Development history: B18--B22 checkpoint and spatial focus}
\label{sec:b18b22}

\subsection{B18 expert-guided epoch selection}
B18 kept the B17 training trajectory unchanged and used the reused expert set only to select one global epoch from five. The epoch-wise selection macro AUCs were 0.618716, 0.665450, 0.651115, 0.639416, and 0.642589, so epoch 2 was selected. The value 0.665450 is a checkpoint-selection statistic, not independent validation performance.

\subsection{B19 joint-focused crop plus vignette}
B19 introduced a 90\% centered crop followed by a cosine vignette. The selected epoch 3 statistic was 0.658131. A same-source Grad-CAM diagnostic showed strong activation along the synthetic vignette boundaries. B19 was therefore rejected as a final preprocessing strategy even though classification remained functional.

\subsection{B20 crop-only focus: active historical model}
B20 retained the 90\% centered crop but removed the vignette entirely. It selected epoch 2 with historical selection statistic 0.667159. A replay under later matched evaluation code produced 0.667407, within the frozen tolerance. B20 is the active historical model because it combines the strongest clean knee-focused formulation with a stable frozen B16 encoder and no synthetic spatial boundary.

The canonical B20 training surface is
\begin{verbatim}
training studies       3120
eligible MRI series   17475
usable B6 cells       14123
positive / negated    6871 / 7252
encoder                frozen B16
selected epoch         2
TTA                    [-1,0,1]
\end{verbatim}

\subsection{B21 pre-resize crop}
B21 corrected the geometric interpretation by applying the 90\% crop before resizing. On the leakage-safe weak-v2 surface B21 improved over its matched B20 control by approximately $+0.01114$, but the single reused-expert acceptance look gave 0.657320 versus B20 replay 0.667407. B21 was not promoted.

\subsection{B22 training-duration audit}
B22 extended the B21 pre-resize-crop trajectory to five epochs. Training loss continued to fall, but the reused expert trajectory peaked at epoch 2 (0.657427) and deteriorated afterward. Longer downstream training therefore did not rescue the pre-resize crop and reinforced the practical fixed-E2 regime used by later B20-family experiments.

\section{Development history: B23--B26 supervision studies}
\label{sec:b23b26}

\subsection{B23 local-LLM report labelling}
B23 replaced the B6 regex parser experimentally with a local open-weights Qwen report labeller while preserving the twelve target meanings and four report states. On the reused expert audit, B23 increased state-only macro AUC from 0.702460 to 0.812516 and increased usable-cell coverage, but specificity fell from 0.6061 to 0.5678. Because improved specificity was part of the predeclared adoption gate, B23-v1 formally failed and no canonical B23 development split was adopted.

\subsection{B24X density pilot}
B24X was explicitly exploratory and did not override the failed B23 gate. On a matched 692-study pilot, the B6 control achieved weak-v2 macro AUC 0.614849, full B23 supervision achieved 0.711613, and a density-preserving policy that kept every B6 cell while filling only B6-silent cells achieved 0.714799. This showed that supervision recovery, rather than replacing existing B6 decisions, explained essentially the entire pilot gain.

\subsection{B25X full hybrid-supervision experiment}
B25X scaled the fill idea to a leakage-safe 2,497-study weak-v2 training surface. The B6 control, pure Hybrid, and B6+Hybrid-fill arms reached 0.672372, 0.726878, and 0.730847 on weak-v2. However, the gain was dominated by Synovitis, where B6 contained only 13 explicit negatives in the B25X training surface and Hybrid-fill recovered 136 additional negatives. Excluding Synovitis, the eleven-target macros were 0.711950 for B6, 0.709133 for Hybrid, and 0.714348 for Fill. The defensible conclusion was therefore class-coverage repair for Synovitis, not broad twelve-target superiority.

\subsection{B26.2 deterministic supervision repair}
The B26 family attempted to turn the supervision-density observation into a deterministic, reproducible repair. B26.2 produced a final Synovitis balance of 475 positive and 112 negative cells, but the exact matched fixed-E2 comparison yielded B20 0.667407 versus B26.2 0.666297, a raw difference of $-0.001109$. Synovitis itself declined materially. Mechanism audit indicated mention-selection bias, so the B26 supervision-repair family was closed and no B26.3 was created.

\section{B27--B33 complementary-series experiments}
\label{sec:b27b33}

\subsection{B27.1 metadata routing}
An initial B27 fluid/fat routing idea was invalid because the relevant metadata were perfectly collinear over the 17,475-series B20 surface. B27.1 instead used plane plus paired-sequence metadata and introduced only 60 routing parameters. It optimized successfully but reached approximately 0.659923, below B20 replay by about 0.00748. The routing family was closed.

\subsection{B28 max-evidence residual}
Let $M$ be the elementwise maximum over the sixteen image-content slice embeddings and $A$ the B20 series token. B28 used
\begin{equation}
T=A+\tanh(g)\odot\LNzero(M),
\end{equation}
with zero-initialized feature-wise gate $g$. At $g=0$ the model is exactly B20. B28 reached approximately 0.63835, a material decline, and the max-evidence residual was closed.

\subsection{B29 learned complementary summary}
B29 added a second softmax summary over the same sixteen B20 slice tokens:
\begin{equation}
C=\LNzero\!\left(\sum_{i=1}^{16}w_iX_i\right),\qquad
w_i=\frac{\exp(q^TX_i/\sqrt{768})}{\sum_j\exp(q^TX_j/\sqrt{768})},
\end{equation}
with learned query $q\in\R^{768}$. The final token is
\begin{equation}
T=A+\tanh(g)\odot(C-A).
\label{eq:b29}
\end{equation}
The query and zero-initialized gate add 1,536 parameters. On the reused expert surface B29 reached 0.676888, raw $+0.009481$ over B20. Its paired median difference was $+0.009421$ with 95\% interval $[-0.003749,+0.024188]$. B29 was frozen as promising but not promoted.

\subsection{B30 projected complementary attention}
B30 reused B20's attention projections in the complementary path while detaching that branch from the shared parameters. Although the mechanism trained, the reused macro fell to approximately 0.654704. The projected-attention formulation was closed.

\subsection{B31 local-context complementary scorer}
B31 retained the B29 complementary query but perturbed its scores using immediate through-plane context:
\begin{equation}
H=X+\mathrm{DWConv}_{k=3}(\LNzero(X)).
\end{equation}
The depthwise Conv1d has 768 groups, kernel size three, no bias, and zero initialization. The complementary scores are computed from $H$, while the weighted values remain the original $X$. B31 adds 3,840 parameters relative to B20.

The mechanism audit showed only a tiny change in the B29 attention distribution: normalized Jensen--Shannon divergence was approximately $3.56\times10^{-9}$, top-1 agreement 98.9\%, and top-3 overlap 99.63\%. Nevertheless, B31 reached 0.682280 on the reused expert surface, the highest B20-family development point estimate. Relative to B20 the raw difference was $+0.014873$; relative to B29 it was $+0.005392$. Because the surface was heavily reused and the mechanism effect was small, B31 was frozen as the leading development candidate rather than promoted.

\subsection{B32 weighted dispersion summary}
B32 tested second-order information with the same complementary attention weights:
\begin{align}
\mu&=\sum_iw_iX_i,\\
\sigma&=\sqrt{\sum_iw_i(X_i-\mu)^2+10^{-6}},
\end{align}
and
\begin{equation}
T=A+\tanh(g_\mu)\odot(\LNzero(\mu)-A)
+\tanh(g_\sigma)\odot\LNzero(\sigma).
\end{equation}
The mean--dispersion residual cosine was approximately 0.074, confirming a non-redundant branch, but B32 reached only 0.668699. Non-redundancy was therefore not sufficient for improved ranking, and the dispersion formulation was closed.

\subsection{B33 exact uniform complementary mean}
B33 removed the learned complementary query entirely and used
\begin{equation}
C_{\mathrm{mean}}=\LNzero\!\left(\frac{1}{16}\sum_{i=1}^{16}X_i\right),
\end{equation}
with
\begin{equation}
T=A+\tanh(g)\odot(C_{\mathrm{mean}}-A).
\label{eq:b33}
\end{equation}
Only the 768-dimensional zero-initialized gate is new. The full-surface E2 mechanism audit showed mean absolute effective gate 0.00502 and maximum 0.02290; the uniform mean and B20 token were nearly orthogonal while the effective injected residual was only about 0.64\% of the B20 token norm.

B33 reached 0.676446 on the reused expert surface, only $-0.000442$ below B29. The paired B33--B29 median was $-0.000496$, 95\% interval $[-0.023312,+0.022036]$. Thus B33 was a successful simplification on the reused surface: it recovered essentially the B29 point estimate using half as many new parameters.

\begin{table}[ht]
\centering
\caption{B20-family reused expert development results. This table is diagnostic only.}
\label{tab:b20family}
\begin{tabular}{llcl}
\toprule
Model & Main intervention & Macro AUC & Status \\
\midrule
B20 & Post-resize 90\% crop-only reference & 0.667407 & Active historical \\
B26.2 & Deterministic supervision repair & 0.666297 & Closed \\
B27.1 & Plane + paired-sequence routing & 0.659923 & Closed \\
B28 & Max-evidence residual & 0.63835 & Closed \\
B29 & Learned complementary summary & 0.676888 & Frozen promising \\
B30 & Projected complementary attention & 0.654704 & Closed \\
B31 & Local-context complementary scorer & 0.682280 & Frozen leading \\
B32 & Weighted dispersion residual & 0.668699 & Closed \\
B33 & Exact uniform complementary mean & 0.676446 & Frozen simplification \\
\bottomrule
\end{tabular}
\end{table}

\section{Prospective weak-validation v1}
\label{sec:pv1}

\subsection{Why a new development surface was required}
By B33, the 58-study expert set had been inspected across many sequential decisions. Continuing architecture development against it would increasingly optimize to the evaluation surface. PV1 was therefore frozen before B34 and the reused expert set was retired as the primary architecture-selection surface.

\subsection{Frozen split}
For StudyInstanceUID $u$, PV1 computes
\begin{equation}
k(u)=\mathrm{SHA256}(\texttt{CNN\_CPC|prospective-weak-v1|2026-08-16}\,\Vert\,\texttt{NUL}\,\Vert\,u).
\end{equation}
All 3,120 B6-active UIDs are sorted deterministically by this key. The first 624 form validation and the remaining 2,496 form training. B6 labels, expert labels, previous model scores, and predictions do not influence membership. The frozen split SHA-256 is
\begin{verbatim}
a0032307abb1ab99724eb39fac25332ce131c575f64d823083bb37f5ec20d1e6
\end{verbatim}
The validation partition contains 624 studies and 3,544 MRI series. The training partition contains 13,931 eligible series and 11,303 usable weak cells, including 5,559 positive and 5,744 negated cells.

The split certifies StudyInstanceUID-level separation only; patient-level separation is not claimed because a trustworthy patient identifier is not part of the frozen contract.

\subsection{Representation-pretraining limitation}
B16 had been aligned using all 4,349 non-gold MRI/report pairs before PV1 existed. Therefore PV1 is valid only for comparing downstream models that share the exact frozen B16 encoder. It may not be used to select a new encoder or representation-pretraining method.

\subsection{Matched B20/B31/B33 training}
The three controls were retrained on exactly the same 2,496-study PV1 training subset with fixed E2 endpoint, frozen B16 encoder, B20 post-resize crop, identical optimizer and schedule, and a matched post-construction stochastic seed. This seed reset prevents extra candidate parameters from shifting the downstream stochastic path merely because model construction consumed a different number of random draws.

\begin{table}[ht]
\centering
\caption{PV1 matched-control training audits. Training loss is descriptive and was not the selection metric.}
\label{tab:pv1train}
\begin{tabular}{lrrrr}
\toprule
Model & New params & E1 loss & E2 loss & Runtime (s) \\
\midrule
B20 & 0 & 0.7667184 & 0.6924914 & 2817.58 \\
B31 & 3840 & 0.7665237 & 0.6912229 & 3196.35 \\
B33 & 768 & 0.7667159 & 0.6833535 & 2981.71 \\
\bottomrule
\end{tabular}
\end{table}

No validation study or expert study entered the gradients, and the frozen encoder fingerprint remained unchanged.

\subsection{Predeclared PV1 metrics}
For target $c$, the primary loss is
\begin{equation}
L_c=\frac{\sum_i w_{ic}\left[-y_{ic}\log p_{ic}-(1-y_{ic})\log(1-p_{ic})\right]}
{\sum_i w_{ic}},
\end{equation}
and the primary selection statistic is
\begin{equation}
L_{\mathrm{PV1}}=\frac{1}{12}\sum_{c=1}^{12}L_c,
\end{equation}
where lower is better. The secondary statistic is macro ROC AUC over B6 positive/negated states for targets where both classes are present. Paired study-level bootstrap comparisons were predeclared for B31--B20, B33--B20, and B33--B31.

\subsection{Memory-safe evaluator}
The first evaluator implementation loaded B20, B31, and B33 simultaneously. Under the local Linux desktop session, \code{systemd-oomd} terminated the GNOME terminal scope after a recorded 56.6~GB memory peak and 3.5~GB swap peak. No scientific result from that failed execution was accepted.

Evaluator v1.1.0 changed resource management only: one checkpoint is loaded at a time, predictions are moved to CPU/disk, the model and payload are deleted, garbage collection and CUDA cache release are performed, and only then is the next checkpoint loaded. Evaluation uses batch size one, one worker, prefetch factor one, no persistent worker, and zero worker-local raw-series cache. TTA remains $[-1,0,1]$, and prediction semantics and metrics are unchanged.

\subsection{Completed PV1 result}
The corrected evaluation completed on all 624 validation studies. Table~\ref{tab:pv1result} gives the prospectively frozen global results.

\begin{table}[ht]
\centering
\caption{Completed PV1 result. Primary loss is lower-is-better. Secondary AUC is computed over hard B6 positive/negated states. PV1 is weak-label architecture-selection evidence, not independent expert validation.}
\label{tab:pv1result}
\begin{tabular}{lcc}
\toprule
Model & Macro weighted soft BCE $\downarrow$ & Macro AUC $\uparrow$ \\
\midrule
B20 & 0.6155808 & 0.5727579 \\
B31 & \textbf{0.5743066} & \textbf{0.7567309} \\
B33 & 0.5849691 & 0.7565223 \\
\bottomrule
\end{tabular}
\end{table}

The primary absolute differences are approximately
\begin{align}
L_{\mathrm{B31}}-L_{\mathrm{B20}}&=-0.04127,\\
L_{\mathrm{B33}}-L_{\mathrm{B20}}&=-0.03061,\\
L_{\mathrm{B33}}-L_{\mathrm{B31}}&=+0.01066.
\end{align}
The paired study bootstrap yielded
\begin{center}
\begin{tabular}{lccc}
\toprule
Comparison & Median difference & 95\% interval & $P$(candidate better) \\
\midrule
B31--B20 & $-0.041141$ & $[-0.051822,-0.029586]$ & 1.0000 \\
B33--B20 & $-0.030688$ & $[-0.047346,-0.013034]$ & 0.9992 \\
B33--B31 & $+0.010862$ & $[+0.002246,+0.019512]$ & 0.0050 \\
\bottomrule
\end{tabular}
\end{center}
Because lower loss is better, all three intervals support the ranking
\begin{equation}
\boxed{\mathrm{B31}\;>\;\mathrm{B33}\;>\;\mathrm{B20}}
\end{equation}
under the predeclared PV1 primary metric.

The secondary AUC gives an additional but different picture. B31 and B33 are almost identical in ranking performance (0.75673 versus 0.75652), while both are far above B20 (0.57276). Therefore the prospectively observed B31--B33 advantage is primarily expressed in the predeclared weighted soft-label loss rather than a meaningful difference in macro AUC. This distinction is important: B31 is the PV1-selected architecture because the primary criterion was frozen in advance, not because every secondary statistic is larger.

\subsection{PV1 decision}
PV1 answers the architecture-selection question clearly. B31 is now the \emph{prospectively selected downstream development architecture} for future fixed-encoder work. B33 remains a successful and informative simplification, but it does not replace B31 under the frozen primary criterion. B20 remains the \emph{active historical model} because PV1 is weak-label validation and does not constitute independent expert or hidden-test evidence.

No target-wise B20/B31/B33 switching, blending, or gate tuning is permitted from the PV1 per-target outputs. The PV1 split must remain unchanged after inspection of these results.

\section{Runtime and execution policy}
\label{sec:runtime}

B20-family training uses one GPU with CPU DICOM/data work. Recent full-surface B29--B33 fixed-E2 runs complete in roughly one hour on the NVIDIA RTX A4500 Laptop GPU used for local experiments. The 80\% PV1 matched-control runs complete in approximately 47--53 minutes each. Competition-oriented workflows retain a conservative internal runtime budget below the nine-hour ceiling.

Resource controls are separated from scientific controls. Batch size, worker count, prefetch depth, and cache size may be reduced to prevent memory failure when MRI sampling, TTA, checkpoint weights, and prediction arithmetic are unchanged. The PV1 v1.1.0 evaluator is an example of this resource-only change.

\section{Reproducibility and artifacts}
\label{sec:reproducibility}

The repository preserves the B0--B33 archive under \code{developments/}. Current top-level code exposes B20 as the canonical historical model, while recent development modules record the frozen encoder fingerprint, series policy, crop policy, supervision contract, new parameter count, gradient checks, coverage, runtime, checkpoint identity, and whether any expert or validation labels entered gradients or checkpoint selection.

PV1 stores the frozen split manifest, B20/B31/B33 matched checkpoints and audits, per-model predictions, paired predictions, and \code{comparison.json}. The completed evaluator reports version 1.1.0, split SHA-256, memory policy, primary and secondary metrics, paired bootstrap comparisons, the frozen encoder fingerprint, and the explicit governance statement that PV1 does not replace independent expert or hidden competition evaluation.

\section{Discussion}
\label{sec:discussion}

The full B0--B33 history reveals several consistent patterns. First, representation quality matters: B13's hierarchical ImageNet protocol produced a major gain over B12, and B16/B17 further established the retained report-aligned frozen encoder. Second, fitting weak supervision more strongly does not guarantee improved expert ranking. B14, B15, and B22 all contain examples where lower training loss or stronger teacher agreement failed to translate into a better reused expert result. Third, data-contract choices are as important as architecture. B23--B26 showed that denser report supervision can repair severe class-coverage failures, particularly Synovitis, while also creating specificity and selection-bias risks.

The B27--B33 sequence narrows the architecture question. Metadata routing, max-evidence pooling, projected complementary attention, and weighted dispersion all failed to justify promotion. B29 demonstrated that a second zero-gated series summary can be useful. B33 then showed on the reused expert surface that a simple exact uniform mean can recover almost the same point estimate as B29. However, the prospective PV1 comparison adds an important new constraint: B31 is better than B33 on the predeclared weighted soft-label loss, while their secondary macro AUCs are almost identical.

This means the B31 mechanism should not be interpreted merely as ``better slice ranking.'' Its local context barely changes the complementary attention distribution, and PV1 shows almost no AUC separation from B33. One plausible interpretation is that B31 provides a more suitable calibrated complementary representation under the B6 soft-label objective. That interpretation remains a hypothesis rather than a causal result. The next mechanistic experiment should therefore be designed globally and prospectively rather than by inspecting target-specific PV1 winners.

PV1 also has clear limitations. Its labels are generated by B6 rather than experts, and the B16 encoder was historically aligned using reports from the complete non-gold population before PV1 existed. Therefore the large B20-to-B31 PV1 improvement establishes out-of-sample downstream agreement with the frozen weak-supervision contract, not clinical superiority. The B31 selection is meaningful for architecture development but does not replace an independent hidden or external expert evaluation.

Additional limitations include sixteen sampled 2.5D positions rather than complete volumetric processing, scanner and physical-scale heterogeneity, no certified patient-level separation in PV1, extreme weak-label imbalance for some targets, and the absence of a current independent hidden-test comparison between B20 and the frozen post-B20 candidates.

\section{Current model status and next experiment}
\label{sec:next}

The current status is deliberately two-tiered:
\begin{equation}
\boxed{\text{B20 = active historical model}},\qquad
\boxed{\text{B31 = PV1-selected development architecture}}.
\end{equation}
B33 remains the preferred simple mechanistic comparator. B29 is also scientifically important because it separates the learned complementary query from B31's added local-context scorer.

Before defining a novel B34 architecture, the most informative low-risk comparison is a matched PV1 retraining of the already frozen B29 architecture. This experiment would not be invented from target-wise PV1 outcomes; B29 predates PV1 and is already frozen. It would decompose the three relevant mechanisms under one prospective surface:
\begin{center}
\begin{tabular}{lll}
\toprule
Model & Complementary summary & Local context \\
\midrule
B33 & exact uniform mean & no \\
B29 & learned query mean-like summary & no \\
B31 & learned query mean-like summary & yes \\
\bottomrule
\end{tabular}
\end{center}
A matched B29 PV1 result would therefore tell whether the B31 advantage over B33 comes mainly from the learned query, from the local-context perturbation, or from their combination. After that single frozen mechanistic comparison, a B34 hypothesis can be defined without returning to the 58-study expert outcomes.

Independent hidden competition or new expert-labelled evaluation remains the required evidence for replacing B20 as the active predictive model.

\section{Conclusions}
\label{sec:conclusion}

CNN\_CPC evolved from early fixed-stream baselines into a leakage-aware hierarchical all-series knee MRI framework with explicit experiment governance. The full development record extends well beyond B9. B10--B17 established the all-series hierarchy and frozen B16 report-aligned encoder; B18--B22 established the active B20 crop-only fixed-E2 regime; B23--B26 exposed both the promise and the risk of denser report supervision; and B27--B33 tested a sequence of increasingly controlled complementary series summaries.

On the reused expert development surface, B31 had the highest B20-family macro AUC, while B33 showed that much of the B29 effect could be reproduced with an exact uniform mean and only 768 new parameters. Because that surface had been reused repeatedly, the project froze PV1 before further architecture development. The completed PV1 comparison is the strongest current architecture-selection evidence: B31 achieves primary macro weighted soft BCE 0.57431 versus 0.58497 for B33 and 0.61558 for B20, with paired bootstrap intervals excluding zero for B31 versus both controls. B31 and B33 have nearly identical secondary macro AUCs, 0.75673 and 0.75652, respectively.

The correct decision is therefore not to promote B31 directly to an independently validated model. Instead, B31 becomes the prospectively selected downstream development architecture, B33 remains the principal simplification comparator, and B20 remains the active historical model pending hidden or external expert evidence. This separation of model development from evaluation reuse is a central contribution of the CNN\_CPC workflow.

\section*{Code availability}
The source code and documentation are available at \url{https://github.com/mtalafha90/CNN_CPC}. Competition data are distributed separately by the competition organizer and are not redistributed by the repository.

\section*{Acknowledgements}
The author acknowledges the University of Sharjah and the Research Institute of Sciences and Engineering for research support and computational facilities. The RSNA Knee Abnormality Detection competition organizers are acknowledged for providing the challenge data and evaluation framework.

\appendix

\section{Exact target order}
\label{app:targets}
\begin{verbatim}
ACL
MCL
Medial Meniscus
Lateral Meniscus
Medial OA
Lateral OA
PF OA
Effusion
Synovitis
Baker's
Contusion
Fracture
\end{verbatim}
Changing this ordering without retraining violates the checkpoint contract.

\section{PV1 audit contract}
\label{app:pv1audit}
\begin{verbatim}
evaluation version      1.1.0
split SHA256            a0032307abb1ab99724eb39fac25332ce131c575f64d823083bb37f5ec20d1e6
training studies        2496
validation studies       624
validation series       3544
training series        13931
training weak cells    11303
positive / negated     5559 / 5744
fixed endpoint         E2
TTA                    [-1,0,1]
validation in gradient 0
expert labels read     false
weak-label validation true
encoder SHA256          b328667cf9dfa9b909ef181c1bcc8975ec42bcd8b9eddad08f908875b73fae96
\end{verbatim}

The frozen low-memory evaluation policy is one resident model, batch size one, one worker, prefetch factor one, no persistent workers, and zero series cache per worker. These controls do not change prediction semantics.

\section{Current experiment decision register}
\label{app:decisionregister}
\begin{longtable}{p{0.12\linewidth}p{0.24\linewidth}p{0.54\linewidth}}
\toprule
Model & Status & Decision \\
\midrule
B20 & Active historical & Canonical checkpoint remains the externally unevaluated working reference. \\
B21/B22 & Closed & Pre-resize crop and longer-duration rescue not supported. \\
B23-v1 & Formal gate failed & Higher coverage/AUC but lower specificity; not adopted. \\
B24X/B25X & Exploratory & Useful supervision-density diagnostics; no promotion. \\
B26.2 & Closed & Deterministic repair did not improve the matched B20 result. \\
B27.1 & Closed & Metadata-routing family not promoted. \\
B28 & Closed & Max-evidence residual declined materially. \\
B29 & Frozen promising & Existing mechanistic candidate; suitable for one matched PV1 decomposition run. \\
B30 & Closed & Projected complementary attention declined. \\
B31 & PV1-selected development architecture & Best predeclared PV1 primary metric; not independent expert validation. \\
B32 & Closed & Dispersion non-redundant but not beneficial. \\
B33 & Frozen simplification & Better than B20 on PV1 but worse than B31 on primary loss; secondary AUC essentially tied with B31. \\
PV1 & Complete/frozen & B31 $>$ B33 $>$ B20 by primary weighted soft-label BCE. \\
B34 & Not yet defined & Define only after prospectively deciding whether to perform the frozen B29 decomposition. \\
\bottomrule
\end{longtable}

\section{Reporting checklist}
Before interpreting a future result, verify the intended checkpoint and endpoint; the B6, crop, series-policy, and encoder fingerprints; the exact evaluation-surface role; aligned StudyInstanceUIDs for paired comparisons; absence of target-wise outcome-derived switching or blending; and whether a resource-only runtime change altered prediction semantics. Promotion claims must be reserved for evidence stronger than the reused 58-study set and PV1 weak-label validation.

\begin{thebibliography}{00}
\bibitem{bien2018mrnet}
N. Bien, P. Rajpurkar, R.L. Ball, J. Irvin, A. Park, E. Jones, M. Bereket, B.N. Patel, K.W. Yeom, K. Shpanskaya, F.G. Halabi, E.J. Zucker, A.M. Fanton, H. Amanatullah, C.F. Beaulieu, G.M. Riley, R.J. Stewart, L. Blankenberg, D. Larson, R.H. Jones, C.P. Langlotz, A.Y. Ng, M.P. Lungren,
Deep-learning-assisted diagnosis for knee magnetic resonance imaging: Development and retrospective validation of MRNet,
PLoS Medicine 15 (2018) e1002699.

\bibitem{liu2022convnext}
Z. Liu, H. Mao, C.-Y. Wu, C. Feichtenhofer, T. Darrell, S. Xie,
A ConvNet for the 2020s,
Proceedings of the IEEE/CVF Conference on Computer Vision and Pattern Recognition (2022) 11976--11986.

\bibitem{vaswani2017attention}
A. Vaswani, N. Shazeer, N. Parmar, J. Uszkoreit, L. Jones, A.N. Gomez, L. Kaiser, I. Polosukhin,
Attention is all you need,
Advances in Neural Information Processing Systems 30 (2017).

\bibitem{paszke2019pytorch}
A. Paszke, S. Gross, F. Massa, A. Lerer, J. Bradbury, G. Chanan, T. Killeen, Z. Lin, N. Gimelshein, L. Antiga, et al.,
PyTorch: An imperative style, high-performance deep learning library,
Advances in Neural Information Processing Systems 32 (2019).
\end{thebibliography}

\end{document}
